\section{Microscopic realizations}
\label{sec:microscopic}

This section verifies the hypotheses of Section~\ref{sec:thresholds}
for the mean-field Potts model, with exact finite-system formulas, and
for the short-range Potts model, via contour estimates. In both, the
reservoir threshold follows from the established phase structure; we
make no new claim about the location or order of the transitions.

\subsection{The three-state Curie--Weiss--Potts model}
\label{subsec:meanfield}

Let $n_1+n_2+n_3=N$ be the occupation numbers of $N$ three-state spins and
take
\begin{equation}
 U_N=-\frac{J}{2N}\sum_{j=1}^3n_j^2,\qquad J>0,
 \qquad \beta J=4\log2.
 \label{eq:mf-hamiltonian}
\end{equation}
Optionally add $2aN$ independent quadratic momentum coordinates, with
$a>0$ fixed and $2aN$ integral along the size sequence. Their canonical
energy $K_N$ is Gamma distributed with shape $aN$ and rate $\beta$, and
$E_N=U_N+K_N$. For $a=0$ there are no momenta and $K_N=0$.

\begin{theorem}[Mean-field microscopic reservoir threshold]
\label{thm:mf-bath}
For every fixed $J>0$ and $a\ge0$ in the model just defined, an admissible
calibration for which $\TV(Q_N,P_N)\to0$ exists if and only if
$c_N\gg N^{3/2}$. The secant calibration~\eqref{eq:secant-calibration}
uses the centers
\begin{equation}
 e_{-,N}=N\left(-\frac J4+\frac a\beta\right),\qquad
 e_{+,N}=N\left(-\frac J6+\frac a\beta\right).
 \label{eq:mf-centers}
\end{equation}
\end{theorem}

The mean-field phase structure and fluctuations are known in greater
generality~\cite{CosteniucEllisTouchette2005,GandolfoRuizWouts2010};
we derive what is needed, making the normalization and the degenerate
disordered spin-energy fluctuation explicit.

Write $p_j=n_j/N$, with coordinates $(p_1,p_2)$ on the probability
simplex and $p_3=1-p_1-p_2$. The rate potential, up to an additive
constant, is
\begin{equation}
 I(p)=\sum_{j=1}^3p_j\log p_j-\frac{\beta J}2\sum_{j=1}^3p_j^2.
 \label{eq:mf-rate}
\end{equation}
Its minima are one disordered and three ordered points,
\begin{equation}
 p^{(d)}=(1/3,1/3,1/3),\qquad
 p^{(o)}\in\{(2/3,1/6,1/6),(1/6,2/3,1/6),(1/6,1/6,2/3)\}.
 \label{eq:mf-minima}
\end{equation}
The list is complete. First, a boundary point cannot minimize $I$:
transferring a small mass $\varepsilon>0$ to a zero coordinate lowers
the entropy term by order $\varepsilon|\log\varepsilon|$ and the
quadratic term by only
$O(\varepsilon)$. At an interior stationary point, all $\log p_j-\beta Jp_j$
are equal. Since $\log x-\beta Jx$ is strictly concave, the coordinates take
at most two values. If two coordinates share the larger value $x$ and
the remaining value is $y<x$, stationarity gives
$\beta J=\log(x/y)/(x-y)>1/x$, and the second variation exchanging the two
large coordinates is negative. A nonsymmetric minimum must therefore
have one large coordinate $r>1/3$ and two equal small coordinates
$(1-r)/2$. Its stationary equation is
\begin{equation}
 \log\frac{2r}{1-r}-\frac{\beta J}2(3r-1)=0.
 \label{eq:mf-scalar-stationarity}
\end{equation}
The derivative of the left side is $1/[r(1-r)]-3\beta J/2$, which has at most
two zeros. By Rolle's theorem, \eqref{eq:mf-scalar-stationarity} has
at most three roots; they are $1/3,1/2,2/3$.
At $r=1/2$ the second variation along this curve is
$4-6\log2<0$. The other candidates have the same value
$I_{\min}=-\log3-(2/3)\log2$, and their Hessians are positive definite:
\begin{align}
 H_d&=(3-\beta J)\begin{pmatrix}2&1\\1&2\end{pmatrix},
 &\det H_d&=3(3-\beta J)^2,
 \label{eq:mf-hessian-d}\\
 H_o&=\begin{pmatrix}15/2-2\beta J&6-\beta J\\6-\beta J&12-2\beta J\end{pmatrix},
 &\det H_o&=3(3-\beta J)(6-\beta J).
 \label{eq:mf-hessian-o}
\end{align}

Partition the simplex into the four Voronoi cells of these minima
(ties resolved deterministically). All four cells have positive limiting
probabilities. Uniform Stirling expansion on compact interior sets gives
the unnormalized occupation weight
\begin{equation}
 \frac{N!}{\prod_j n_j!}e^{-\beta U_N(n)}
 =\frac{1+o(1)}{2\pi N\sqrt{p_1p_2p_3}}\,e^{-NI(p)}.
 \label{eq:mf-stirling}
\end{equation}
Near a minimum, the rescaled two-dimensional lattice has mesh
$N^{-1/2}$ and hence density $N$. Taylor expansion of $I$
in~\eqref{eq:mf-stirling} and Gaussian Riemann sums give a phase
contribution proportional to
$(p_1p_2p_3)^{-1/2}(\det H_i)^{-1/2}$ and a conditional occupation
limit $\Normal(0,H_i^{-1})$; the uniform bound proved below controls
the tails. With minus denoting the combined ordered phases, the phase
weights are
\begin{equation}
 r_*:=\sqrt{\frac{2(3-\beta J)}{6-\beta J}},\qquad
 w_-:=\frac{3r_*}{1+3r_*},\qquad w_+:=\frac1{1+3r_*}.
 \label{eq:mf-phase-weights}
\end{equation}
The tangent gradient of $u(p)=-J\sum_jp_j^2/2$ is
$(-J/2,0)$ at the ordered minimum $p^{(o)}=(2/3,1/6,1/6)$ and zero at
the disordered one. The delta method and independence of the kinetic
energy give the conditional weak energy limits
\begin{equation}
 \frac{E_N-e_{i,N}}{\sqrt N}\ \Rightarrow\ \Normal(0,v_i),
 \qquad
 v_-:=\frac a{\beta^2}+\frac{J^2}{6(3-\beta J)},\quad
 v_+:=\frac a{\beta^2}.
 \label{eq:mf-variances}
\end{equation}
For $a=0$, $\Normal(0,0)$ is the point mass at zero: the disordered
potential energy fluctuates only on scale one, since its first
derivative on the simplex vanishes.

\begin{lemma}[Uniform positive-phase energy moments]
\label{lem:mf-exponential}
For each fixed $a\ge0$, there are $\eta>0$ and $M<\infty$ such that
the two positive energy phase laws defined by the Voronoi partition
satisfy
\begin{equation}
 \sup_{N\ge N_0,i\in\{-,+\}}
 \E\left[\exp\left(\eta\frac{|E_N-e_{i,N}|}{\sqrt N}\right)
       \,\middle|\,i\right]\le M.
 \label{eq:mf-moment-bound}
\end{equation}
Here $N_0$ is fixed and large enough that every conditioning phase has
positive probability.
\end{lemma}
\begin{proof}
We first prove a Gaussian bound for the occupation law. Let
$\mathcal M$ be the set of the four minima. Compactness, the complete
classification of minima, and their positive definite Hessians give
$I(p)-I_{\min}\ge c_0\operatorname{dist}(p,\mathcal M)^2$ globally.
The Gaussian neighborhoods in~\eqref{eq:mf-stirling} give a partition
function lower bound $c e^{-NI_{\min}}$. On a fixed interior subset,
Stirling bounds then give
\[
 P_N(n)\le\frac CN e^{-cN\operatorname{dist}(n/N,\mathcal M)^2}.
\]
It extends globally with a smaller $c$: on a boundary strip disjoint
from $\mathcal M$, where the rate gap is at least $\zeta>0$, the
multinomial type bound
$N!/\prod n_j!\le e^{-N\sum p_j\log p_j}$ gives
$P_N(n)\le C e^{-N[I(p)-I_{\min}]}$ there. Splitting the exponent
between the gap and distance bounds recovers the missing $1/N$, since
$N e^{-N\zeta/2}$ is bounded.

In the cell of $p^{(i)}$, the distance to $\mathcal M$ is
$|p-p^{(i)}|$ and the Lipschitz property of $u$ gives
$|U_N-Nu(p^{(i)})|/\sqrt N\le C\sqrt N|p-p^{(i)}|$; completing
the square in the lattice Gaussian sum shows
\[
 \frac CN\sum_{n_1,n_2}
 e^{-cN|n/N-p^{(i)}|^2+\eta C\sqrt N|n/N-p^{(i)}|}\le C_\eta.
\]
Conditioning preserves it, as the phase probabilities are bounded
below for large $N$. For $a>0$, the centered Gamma transform is
\[
 \E e^{t(K_N-aN/\beta)/\sqrt N}
 =e^{-ta\sqrt N/\beta}
   \left(1-\frac{t}{\beta\sqrt N}\right)^{-aN}.
\]
It is uniformly bounded for fixed $t$ and large $N$. Using
$e^{t|x|}\le e^{tx}+e^{-tx}$ and independence completes the proof.
\end{proof}

\begin{proof}[Proof of Theorem~\ref{thm:mf-bath}]
The gap of the two energy centers is $NJ/12$.
Lemma~\ref{lem:mf-exponential} supplies a positive decomposition with
no exceptional component and $\sqrt N$ moment bounds, and
\eqref{eq:mf-variances} a nondegenerate ordered weak limit
($v_->0$ because $\beta J<3$).
Apply Corollary~\ref{cor:two-phase-iff}.
\end{proof}

\subsection{Continuous energy limits and exact finite-system formulas}
\label{subsec:mf-exact}

For $a>0$ the mean-field law satisfies the density hypotheses of
Section~\ref{subsec:density-tails}; this subsection proves this and
derives the exact finite-$N$ formulas used in
Section~\ref{sec:numerics}.

For $a>0$, the standardized Gamma density converges uniformly to the
Gaussian density of variance $a/\beta^2$, by Stirling's formula on
compact standardized intervals and, outside them, unimodality: once the
mode is inside, the supremum outside is at the nearest endpoint, and
enlarging the interval bounds both tails uniformly.
Convolution with the conditional spin-energy weak limits (the $a=0$
case of~\eqref{eq:mf-variances}) therefore gives
the local $L^1$ density limits~\eqref{eq:local-density-assumption} with
weights~\eqref{eq:mf-phase-weights}, and the stronger interior bound
\begin{equation}
 p_N(E)\le\frac C{\sqrt N}\sum_{i\in\{-,+\}}
          \exp\left[-c\frac{(E-e_{i,N})^2}{N}\right]
 \quad(e_{-,N}\le E\le e_{+,N}).
 \label{eq:mf-density-envelope}
\end{equation}
To prove it, note that the kinetic energies allowed in this interval
obey $0<K\le C_0N$, where the kinetic density $g_N$ has logarithmic
curvature
$-(aN-1)/K^2\le-c_1/N$ for all large $N$, and its maximum is
$O(N^{-1/2})$. Expanding at its mode, which differs from $aN/\beta$
by a constant, gives
$g_N(K)\le C N^{-1/2}e^{-c(K-aN/\beta)^2/N}$ on this interval.
For a nearest minimum $p^{(i)}$, Lipschitz continuity of $u$ yields
\[
 |p-p^{(i)}|^2+|E/N-u(p)-a/\beta|^2
 \ge c\{|p-p^{(i)}|^2+|E/N-u(p^{(i)})-a/\beta|^2\}.
\]
Multiply the occupation bound from Lemma~\ref{lem:mf-exponential} by
the Gamma bound and sum over each cell; the lattice Gaussian sum is
$O(N)$ and cancels the prefactor $N^{-1}$, proving
\eqref{eq:mf-density-envelope}. This also controls contamination from
the other phase in the local density limit.

For finite systems set $A=aN>0$ and $m(n)=N!/(n_1!n_2!n_3!)$. The
canonical and reservoir occupation laws are exactly
\begin{align}
 P_N(n)&=\frac{m(n)e^{-\beta U_N(n)}}
                  {\sum_{n'}m(n')e^{-\beta U_N(n')}},
 \label{eq:mf-canonical-occupancy}\\
 Q_N(n)&=\frac{m(n)[\mathcal E_N-U_N(n)]_+^{A+c_N}}
                  {\sum_{n'}m(n')[\mathcal E_N-U_N(n')]_+^{A+c_N}}.
 \label{eq:mf-bath-occupancy}
\end{align}
Conditional on $n$, $K_N$ is canonical Gamma$(A,\beta)$, whereas
\begin{equation}
 \frac{K_N}{\mathcal E_N-U_N(n)}\ \bigg|\ n
 \sim\operatorname{Beta}(A,c_N+1)
 \quad\hbox{under }Q_N.
 \label{eq:mf-beta-conditional}
\end{equation}
Equations~\eqref{eq:mf-bath-occupancy}
and~\eqref{eq:mf-beta-conditional} follow from integrating the product
of subsystem and reservoir densities:
\[
 \int_0^{\mathcal E_N-U_N(n)}
 K^{A-1}[\mathcal E_N-U_N(n)-K]^{c_N}\dd K
 =\mathrm B(A,c_N+1)[\mathcal E_N-U_N(n)]^{A+c_N}.
\]
For $a=0$ the occupation formulas hold with $A=0$ and $K_N=0$; no
shape-zero Gamma or Beta law is used.

Let $G_A$ be the CDF of a unit-rate Gamma law and $I_x(A,c_N+1)$ the
regularized incomplete Beta function, extended by zero for $x\le0$
and one for $x\ge1$. Then, for $a>0$,
\begin{align}
 P_N(E_N\le t)&=\sum_nP_N(n)G_A\bigl(\beta[t-U_N(n)]\bigr),
 \label{eq:mf-gamma-cdf}\\
 Q_N(E_N\le t)&=\sum_{n:Q_N(n)>0}Q_N(n)
 I_{[t-U_N(n)]/[\mathcal E_N-U_N(n)]}(A,c_N+1),
 \label{eq:mf-beta-cdf}
\end{align}
where $G_A(x)=0$ for $x\le0$. These are exact finite sums. The log
likelihood in~\eqref{eq:exact-marginal} is
strictly concave below its cutoff and tends to $-\infty$ at both ends
of that half-line; for $a>0$ it is not constant on the canonical
support, so its normalized maximum exceeds zero. The normalized log
likelihood is therefore positive exactly on an interval $(t_L,t_R)$
between its two roots, and
\begin{equation}
 \TV(Q_N,P_N)=Q_N(t_L<E_N<t_R)-P_N(t_L<E_N<t_R).
 \label{eq:mf-tv-cdf}
\end{equation}
Equations~\eqref{eq:mf-gamma-cdf}--\eqref{eq:mf-beta-cdf} evaluate
this directly, even when an endpoint lies outside the subsystem
support. For the pure-spin model one instead sums
the positive likelihood differences over the discrete occupation law.

\subsection{A short-range Potts theorem}
\label{subsec:short-range}

Fix a dimension $d\ge2$ and an integer $q\ge q_0(d)$, independent of
system size; the threshold $q_0(d)>2d$ is fixed at the start of
Appendix~\ref{app:short-range-proof}, large enough for the cited
contour results. On the periodic torus
$\Lambda_L=(\mathbb Z/L\mathbb Z)^d$, with $N=L^d$ and the set
$\mathcal B_L$ of its $dN$ nearest-neighbor bonds, let
\begin{equation}
 U_N(\sigma)=-\sum_{\{x,y\}\in\mathcal B_L}
                   \ind_{\{\sigma_x=\sigma_y\}},
 \qquad \sigma\in\{1,\ldots,q\}^{\Lambda_L}.
 \label{eq:sr-hamiltonian}
\end{equation}
The unit coupling fixes energy units, and $\beta_c$ denotes the
coexistence inverse temperature.

\begin{theorem}[Short-range microscopic reservoir threshold]
\label{thm:sr-bath}
For every fixed $d\ge2$ and fixed $q\ge q_0(d)$, the canonical
law of~\eqref{eq:sr-hamiltonian} at $\beta_c$ can be reproduced in
total variation by an exact reservoir~\eqref{eq:bath-dos}, with some
calibration, if and only if $c_N\gg N^{3/2}$. The sufficient
choice is~\eqref{eq:secant-calibration} with centers $Nu_o$ and $Nu_d$,
where $u_o<u_d$ are the two coexistence energy densities.
\end{theorem}

In contrast to the boundary results of
Section~\ref{sec:boundary-smooth}, this threshold holds without
restriction in $d=2$, for the plain spin model. Its proof in
Appendix~\ref{app:short-range-proof} uses two positive phase
partitions that need not coincide: midpoint energy events for
necessity, through Theorem~\ref{thm:two-scale-necessity} with gap
$N(u_d-u_o)$ and $s_N=\sqrt N$, and contour events in the
Edwards--Sokal joint space of spins and bonds for sufficiency, through
Theorem~\ref{thm:positive-sufficiency} and projection to the spins.
Contour estimates control $\sqrt N$-scale exponential moments of the
occupied-bond count, but sufficiency needs them for the spin energy.
The conditional binomial law of bonds given spins provides this
transfer; differentiating a bond-conditioned spin law would not justify
it. The exceptional contour mass decays on the surface scale
$L^{d-1}$, which dominates the largest reservoir gain when
$c_N\gg N^{3/2}$.
